% Inclusion fragment. Requires amsmath, amssymb, amsthm and theorem,
% proposition/corollary environments. No private macros are used.

\subsection{Counter congruence alone does not recover root replacement}
\label{rootind:section}

The two restricted equality rules are not interchangeable. The following
counterinterpretation shows that simply deleting root replacement loses
even the theorem $0=0$. This is a statement about the explicitly specified
native inference system. It does not rule out every different rule,
encoding, numerical basis, or interpretation.

\begin{theorem}[Root replacement is not derivable from the remaining rules]
\label{rootind:theorem}
Let $\mathcal T^{-\mathsf{RL}}$ have the same terms, Eq/$\mathsf D$
object formulas, internal $\mathsf Q$ clauses, and rules R0, RS,
$\mathsf{Hyp}$, $\mathsf{CC}$, $\mathsf{QE}$, $\mathsf{DE}$,
$\mathsf{DI}$-case, and $\mathsf{DI}$-progress as the counter-congruence
calculus, but omit $\mathsf{RL}$. Do not add an unlisted reflexivity,
equality-substitution, or other ordinary equality rule. Then
\[
                    \varnothing\nvdash_{\mathcal T^{-\mathsf{RL}}}0=0.
\]
In particular, $\mathsf{RL}$ is not admissible in this reduced system.
\end{theorem}
\begin{proof}
For this syntactic nonderivability argument, interpret the native atom
written $s=t$ by a binary predicate $E$. We are not assuming that this
predicate is actual equality: the missing equality principles are
precisely what the proposed reduction would have to derive.

Take numerical domain $\{0,1\}$. Interpret zero as $0$ and every
recursion expression by
\[
       \llbracket R(n,a;u,v.h)\rrbracket_\rho=1
\]
under every assignment $\rho$, independently of its counter, seed,
update, and passive parameters. This is compatible with capture-free
substitution and alpha-renaming. The successor abbreviation
$Sa=R(0,a;u,v.0)$ is therefore the constant function $1$, and the
numeral $1=S0$ denotes the domain element $1$.
Interpret the equality-shaped atom by
\[
       \llbracket s=t\rrbracket_\rho\text{ is true}
       \quad\Longleftrightarrow\quad
       \llbracket s\rrbracket_\rho=1.
 \tag{Predicate interpretation}
\]
Its right-hand term is immaterial. The equality on the right of this
metatheoretic definition is the usual equality of the two domain elements.
Interpret the binders and internal clauses by their ordinary meanings:
\[
 \begin{aligned}
  (\mathsf Q_y B)^\circ
    &:=B^\circ(0)\to\forall y\,B^\circ(y),\\
  (\mathsf D_{x,y}B)^\circ
    &:=\exists x\bigl(B^\circ(x,0)\to\forall y\,B^\circ(x,y)\bigr).
 \end{aligned}
\]
These are ordinary two-valued semantics, equivalently a one-world
intuitionistic interpretation.

We check every retained primitive rule under arbitrary assignments and
arbitrary true retained contexts.

\paragraph{Computation and counter congruence.}
Each R0 or RS axiom has an $R$ expression as the left side of its atom.
That expression has value $1$, so the atom is true, irrespective of its
right side. A $\mathsf{CC}$ conclusion also has an $R$ expression on
its left and is therefore true. These rules are sound even though the
binary predicate is not reflexive at zero.

\paragraph{Hypotheses, availability, and elimination.}
An object hypothesis is true whenever its retained context is true.
An available major is either such a retained assumption or an already
proved object formula, so availability preserves truth.
For $\mathsf{QE}$, a true clause
$B(0)\to\forall yB(y)$ and a true base $B(0)$ give its requested
instance $B(t)$. For $\mathsf{DE}$, the true existential major supplies
a domain element satisfying the retained closure clause. The fresh
witness parameter can be assigned that element; its absence from the
outer conclusion and retained context gives the ordinary soundness
argument for existential elimination.

\paragraph{Case introduction for arbitrary object matrices.}
Fix an assignment making $\Gamma$ true, and let $b$ be the value of
the submitted witness $t$. Freshness makes $b$ and $\Gamma$
independent of the new successor parameter $z$. A sound premise
\[
        \Gamma,B(t,0)\vdash B(t,Sz)
\]
ensures that $B(b,0)$ implies $B(b,1)$, since $Sz$ always denotes
$1$. The domain has just $0$ and $1$, so this implication ensures
$B(b,0)\to\forall yB(b,y)$. Choosing $x=b$ proves the intended
$\mathsf D$ conclusion. This verifies $\mathsf{DI}$-case without
restricting its matrix to an equation.

\paragraph{Progress introduction for literal equations.}
Fix the passive parameters and the witness. By inspection of finite
numerical term syntax, a term $U(y)$ has one of the following three
value functions of the remaining coordinate $y$:
\[
                y\longmapsto0,\qquad
                y\longmapsto1,\qquad
                y\longmapsto y.
\]
Indeed a zero term is constant $0$; a fixed passive variable is constant
$0$ or $1$; the distinguished variable has value $y$; and any term
with an $R$ root, including a successor abbreviation, is constant $1$.
For a literal equation $E(y)=(U(y)=V(y))$, its truth depends only on
$U(y)$ by the predicate interpretation. Its truth set is consequently
one of
\[
                       \varnothing,\qquad\{1\},\qquad\{0,1\}.
\]
In every case, truth of $E(0)$ implies truth of $E(y)$ for both domain
elements. Thus the required closure
$E(0)\to\forall yE(y)$ is valid for every fixed witness, even before
using the progress premise. This proves soundness of
$\mathsf{DI}$-progress. The restriction to a literal equation is
essential to this verification; it is not a claim about arbitrary
compound matrices.

All primitive rules of $\mathcal T^{-\mathsf{RL}}$ are now sound.
Induction on a finite derivation therefore preserves truth in this
interpretation. But the atom $0=0$ is false because its left term has
value $0$. It has no derivation from the empty context.

Finally, in the full calculus RS gives
\[
                  R(S0,0;u,v.u)=0.
\]
Two copies of this computation, used as the major and minor of
$\mathsf{RL}$, prove $0=0$. If $\mathsf{RL}$ were admissible in
the reduced system, this three-node argument could be translated into
it, contradicting the counterinterpretation.
\end{proof}

The conclusion is deliberately limited. Deleting $\mathsf{RL}$ destroys
the established identity-language equivalence with the unrestricted
transport calculus. The counterinterpretation is not asserted to be a
model of an ordinary theory in which equality already carries its usual
logical axioms, and it is not an impossibility theorem for all
bi-interpretations with a different equality or domain. Importing those
usual equality axioms silently would change the proof-checker problem.
